\documentclass[11pt]{article}
\usepackage[a4paper,margin=1in]{geometry}
\usepackage{amsmath,amssymb,amsfonts}
\usepackage{mathtools}
\usepackage{hyperref}
\usepackage{enumitem}
\hypersetup{colorlinks=true,linkcolor=blue,citecolor=blue,urlcolor=blue}
\setlist{nosep}
\numberwithin{equation}{section}

\title{Unified Boundary Field Theory of Interest Rates III: The Stochastic Long Boundary \\\large Volatility-Aligned Geometry, Affine Admissibility, and Two-Channel MSNR for Instantaneous Rolling Forward Rates}

\author{Sachin Oza}
\date{\today}

\begin{document}
\maketitle

\begin{abstract}
This paper completes the theoretical construction of the Unified Boundary Field Theory for instantaneous rolling forward rates (IRFRs). Earlier parts of the theory show that a one-boundary forward-rate system admits non-unique affine drift and volatility structures with constant market signal-to-noise ratios. The present paper extends that construction by adding a stochastic long-boundary and deriving the corresponding two-boundary geometry.

The central object is the maturity-loading relation
\[
x_{t,\tau+t}-X_t
=
\left(x_{t,t}-X_t\right)e^{-a^2\tau},
\]
which expresses the remaining-maturity \(\tau\) instantaneous rolling forward rate, at calendar time \(t\), as an exponentially weighted interpolation between the short-boundary state and the stochastic long-boundary. This relation provides both an expected shape condition and a local volatility-alignment condition.

The paper then derives the two-variable Kolmogorov forward equation for the short-boundary and long-boundary coordinates, transforms it into the completed boundary coordinate system, and matches coefficients. The coefficient matching shows that the long-boundary extension does not destroy the affine admissible class. Instead, admissibility requires the coordinate geometry to be aligned with the transformed volatility structure. The diffusion equations also imply that the original and transformed Brownian correlations cannot be chosen independently; on the orientation-preserving branch they coincide as a single structural correlation parameter.

Under affine volatility functions, constant drift-to-volatility ratios, and constant correlation, the completed two-boundary system preserves a global constant market signal-to-noise ratio (MSNR). The result provides the analytic bridge between the one-boundary admissibility theory and the empirical two-boundary covariance structure studied in the companion empirical work.
\end{abstract}

\section{Introduction}

The instantaneous rolling forward rate (IRFR) boundary-field framework begins from the premise that the term structure is a maturity field constrained by boundary conditions. Paper 1 developed the basic maturity-field representation and the exponential boundary transport clock. Paper 2 classified short end boundary regimes and showed that the Robin or mixed short end condition generates the repeated root loading
\[
 \tau e^{-a^2\tau}.
\]
This paper completes the theoretical construction by adding a stochastic realised long-maturity boundary.

The fixed long-boundary theory identifies the asymptotic state with an equilibrium anchor \(x_\infty\). That is useful as a long-run reference, but it is too restrictive when the realised long end of the curve can move persistently. The present paper therefore separates the equilibrium anchor \(x_\infty\) from the realised long-boundary state \(X_t\). The long-boundary displacement is
\begin{equation}
\label{eq:Y-definition}
Y_t:=X_t-x_\infty.
\end{equation}

The completed maturity geometry is generated by three loadings:
\begin{equation}
\label{eq:three-loadings}
e^{-a^2\tau},
\qquad
\tau e^{-a^2\tau},
\qquad
1-e^{-a^2\tau}.
\end{equation}
\begin{itemize}
    \item The first loading transports the realised short boundary into maturity space. 
    \item The second is the Robin interior deformation inherited from the short-boundary classification. 
    \item The third transports the stochastic long-boundary into maturity space.
\end{itemize}

In the canonical bridge representation the curve is written
\begin{equation}
\label{eq:canonical-bridge-intro}
x_t(\tau)
=
x_{t,t}e^{-a^2\tau}
+
B_tx_{t,t}\tau e^{-a^2\tau}
+
X_t\left(1-e^{-a^2\tau}\right).
\end{equation}
The corresponding maturity-space Jacobian is
\begin{equation}
\label{eq:intro-jacobian}
J(\tau)
=
\left(
 e^{-a^2\tau},
 \tau e^{-a^2\tau},
 1-e^{-a^2\tau}
\right).
\end{equation}
The first two components are inherited from the short end boundary theory. The third is the new stochastic long-boundary component.

The main theoretical question is whether this added long-boundary destroys the affine admissible family of the one boundary construction and the answer is no. The long-boundary imposes a compatibility condition on the coordinate geometry, but it does not remove the affine admissible class. The required compatibility is volatility alignment:
\begin{equation}
\label{eq:intro-vol-align}
y_x=\frac{g_y}{g_x},
\qquad
Y_X=\frac{g_Y}{g_X}.
\end{equation}
Once this is imposed, the remaining diffusion equations identify the compatible correlation branches \(\rho_{yY}=\pm\rho_{xX}\). On the orientation-preserving branch, affine \(f\) and \(g\) functions, constant \(f/g\), and a single constant correlation preserve a global two-boundary market signal to noise ratio (MSNR).
The paper proceeds as follows:

\begin {itemize}
 \item Section 2 states the stochastic long-boundary decomposition. \item Section 3 records the instantaneous expected level and full instantaneous volatility implied by the exponential maturity loading. \item Section 4 derives the two-variable Kolmogorov forward equation. \item Section 5 performs the coordinate transform and coefficient matching. \item Section 6 proves affine admissibility and the global correlated MSNR. \item Section 7 distinguishes the Robin covariance-norm closure from the genuine two-channel Robin MSNR. \item Section 8 records the asymptotic boundary dynamics and long-end variance floor. \item Section 9 discusses empirical implications.
\end {itemize}


\section{Stochastic long-boundary decomposition}

Let \(\tau=T-t\) denote remaining maturity. In the fixed long-boundary theory, the asymptotic boundary is represented by an equilibrium anchor \(x_\infty\). The stochastic long-boundary extension separates this equilibrium anchor from the realised long-boundary state \(X_t\). We write the realised displacement from equilibrium as
\[
Y_t:=X_t-x_\infty.
\]

The long-boundary loading is
\[
1-e^{-a^2\tau}.
\]
It vanishes at the short end and tends to one at the long end. Thus the stochastic long-boundary contribution is
\[
\left(1-e^{-a^2\tau}\right)(X_t-x_\infty).
\]

For a short-end boundary branch \(j\), the maturity field may be decomposed as
\begin{equation}
\label{eq:boundary-decomposition}
x_t^j(\tau)
=
x_\infty\left(1-e^{-a^2\tau}\right)
+
\left(1-e^{-a^2\tau}\right)(X_t-x_\infty)
+
Z_t^j(\tau).
\end{equation}
The first term is the fixed-anchor transport component. The second term transports the realised long-boundary displacement into maturity space. The residual \(Z_t^j(\tau)\) is the finite-maturity shape component associated with the short-end boundary branch.

Equivalently, since \(Y_t=X_t-x_\infty\),
\[
x_t^j(\tau)
=
x_\infty\left(1-e^{-a^2\tau}\right)
+
\left(1-e^{-a^2\tau}\right)Y_t
+
Z_t^j(\tau).
\]
The split form is retained because it distinguishes the equilibrium anchor from the realised long-boundary displacement, even though the first two terms combine algebraically to
\[
X_t\left(1-e^{-a^2\tau}\right).
\]

At the short end,
\[
1-e^{-a^2\cdot 0}=0,
\]
so the stochastic long boundary does not move the instantaneous endpoint. At the long end,
\[
1-e^{-a^2\tau}\to1
\qquad
\text{as}
\qquad
\tau\to\infty.
\]
Therefore, if the finite-maturity residual satisfies
\[
Z_t^j(\tau)\to0,
\]
then
\[
\lim_{\tau\to\infty}x_t^j(\tau)=X_t.
\]
The stochastic long-boundary extension therefore replaces the fixed asymptotic limit \(x_\infty\) with the realised long-boundary state \(X_t\), while preserving \(x_\infty\) as the equilibrium reference point.

For the Robin branch, the canonical bridge is
\begin{equation}
\label{eq:canonical-bridge}
x_t(\tau)
=
x_{t,t}e^{-a_R^2\tau}
+
B_tx_{t,t}\tau e^{-a_R^2\tau}
+
X_t\left(1-e^{-a_R^2\tau}\right).
\end{equation}
The short endpoint is pinned to \(x_{t,t}\), the long endpoint is pinned to \(X_t\), and the interior deformation is carried by the repeated-root Robin loading
\[
\tau e^{-a_R^2\tau},
\]
which vanishes at both endpoints. The parameter \(a_R\) is written separately to allow the Robin maturity clock to differ from the corresponding free or Neumann clock.

\section{Instantaneous expected level and volatility loading}

The original maturity-loading intuition is most transparent when \(T_x=t\). Then the short-side coordinate is \(x_{t,t}\), the maturity-\(T\) coordinate is \(x_{t,T}\), and the long-boundary state is \(X_t\). The exponential boundary relation is
\begin{equation}
\label{eq:expected-level-boundary-relation}
x_{t,T}-X_t
=
\left(x_{t,t}-X_t\right)e^{-a^2(T-t)}.
\end{equation}
Equivalently,
\begin{equation}
\label{eq:expected-level-direct}
x_{t,T}
=
e^{-a^2(T-t)}x_{t,t}
+
\left(1-e^{-a^2(T-t)}\right)X_t.
\end{equation}
Therefore the instantaneous conditional expected level at maturity \(T\), given \(x_{t,t}\) and \(X_t\), is
\begin{equation}
\label{eq:instantaneous-expected-level}
\mathbb E_t[x_{t,T}]
=
e^{-a^2(T-t)}x_{t,t}
+
\left(1-e^{-a^2(T-t)}\right)X_t.
\end{equation}
Equivalently,
\[
\mathbb E_t[x_{t,T}-X_t]
=
\left(x_{t,t}-X_t\right)e^{-a^2(T-t)}.
\]

Differentiating \eqref{eq:expected-level-direct} with respect to \(x_{t,t}\), holding \(X_t\) fixed, gives
\begin{equation}
\label{eq:local-maturity-loading-derivative}
\frac{\partial x_{t,T}}{\partial x_{t,t}}
=
e^{-a^2(T-t)}.
\end{equation}
The coefficient-matching result below requires local volatility alignment:
\[
\frac{\partial x_{t,T}}{\partial x_{t,t}}
=
\frac{g_{t,T}}{g_{t,t}}.
\]
Therefore
\begin{equation}
\label{eq:maturity-vol-alignment}
\frac{g_{t,T}}{g_{t,t}}
=
e^{-a^2(T-t)},
\qquad
\text{or}
\qquad
g_{t,T}=e^{-a^2(T-t)}g_{t,t}.
\end{equation}
This is the local maturity-loading volatility relation when the long-boundary is held fixed.

If the long-boundary is itself stochastic, the full instantaneous volatility of \(x_{t,T}\) also contains boundary risk. Suppose
\begin{equation}
\label{eq:short-sde}
dx_{t,t}=-f_t\,dt+\sqrt{2}\,g_t\,dW_t^t,
\end{equation}
and
\begin{equation}
\label{eq:X-sde-simple}
dX_t=-f_X\,dt+\sqrt{2}\,g_X\,dW_t^X,
\end{equation}
with
\begin{equation}
\label{eq:rho-simple}
dW_t^t\,dW_t^X=\rho\,dt.
\end{equation}
Over an instantaneous step, treating the maturity-loading coefficient as fixed, \eqref{eq:expected-level-direct} gives
\[
dx_{t,T}
=
e^{-a^2(T-t)}dx_{t,t}
+
\left(1-e^{-a^2(T-t)}\right)dX_t.
\]
Substitution of \eqref{eq:short-sde} and \eqref{eq:X-sde-simple} yields
\begin{align}
\label{eq:maturity-T-sde}
dx_{t,T}
&=
-\left[
e^{-a^2(T-t)}f_t
+
\left(1-e^{-a^2(T-t)}\right)f_X
\right]dt
\nonumber\\
&\quad+
\sqrt{2}\left[
e^{-a^2(T-t)}g_t\,dW_t^t
+
\left(1-e^{-a^2(T-t)}\right)g_X\,dW_t^X
\right].
\end{align}
Hence the instantaneous drift coefficient is
\begin{equation}
\label{eq:instantaneous-drift-maturity-T}
f_{t,T}
=
e^{-a^2(T-t)}f_t
+
\left(1-e^{-a^2(T-t)}\right)f_X.
\end{equation}
The full two-boundary instantaneous variance coefficient is
\begin{align}
\label{eq:instantaneous-vol-maturity-T}
g_{t,T}^2
&=
e^{-2a^2(T-t)}g_t^2
+
\left(1-e^{-a^2(T-t)}\right)^2g_X^2
\nonumber\\
&\quad+
2\rho e^{-a^2(T-t)}\left(1-e^{-a^2(T-t)}\right)g_tg_X.
\end{align}
Thus the exponential maturity loading governs both the instantaneous expected level and the local volatility alignment, while the full stochastic-boundary volatility is a correlated combination of short-side and long-boundary risk.

\section{Two variable KFE before transformation}

We now record the two variable Kolmogorov forward equation that underlies the coefficient matching argument. Let the original coordinates be \((x,X)\), where \(x\) is the short-side or interior state and \(X\) is the long-boundary state. Use the monograph convention
\begin{equation}
\label{eq:xX-sde}
dx_t=-f_x\,dt+\sqrt{2}\,g_x\,dW_t^x,
\qquad
dX_t=-f_X\,dt+\sqrt{2}\,g_X\,dW_t^X,
\end{equation}
with
\begin{equation}
\label{eq:xX-corr}
dW_t^x\,dW_t^X=\rho_{xX}\,dt.
\end{equation}
For a density \(\varphi=\varphi(x,X,t)\), the KFE is
\begin{align}
\label{eq:xX-KFE-conservative}
\frac{\partial\varphi}{\partial t}
&=
\frac{\partial}{\partial x}(f_x\varphi)
+
\frac{\partial}{\partial X}(f_X\varphi)
+
\frac{\partial^2}{\partial x^2}(g_x^2\varphi)
\nonumber\\
&\quad+
2\rho_{xX}\frac{\partial^2}{\partial x\,\partial X}(g_xg_X\varphi)
+
\frac{\partial^2}{\partial X^2}(g_X^2\varphi).
\end{align}
Expanding and collecting derivatives of \(\varphi\), write
\begin{equation}
\label{eq:xX-KFE-expanded}
\frac{\partial\varphi}{\partial t}
=
A_0\varphi+A_x\varphi_x+A_X\varphi_X+A_{xx}\varphi_{xx}+A_{xX}\varphi_{xX}+A_{XX}\varphi_{XX},
\end{equation}
where
\begin{align}
\label{eq:A0}
A_0
&=
\frac{\partial f_x}{\partial x}
+
\frac{\partial f_X}{\partial X}
+
\frac{\partial^2 g_x^2}{\partial x^2}
+
2\rho_{xX}\frac{\partial^2(g_xg_X)}{\partial x\,\partial X}
+
\frac{\partial^2 g_X^2}{\partial X^2},
\end{align}
\begin{align}
\label{eq:AxAX}
A_x
&=
f_x+2\frac{\partial g_x^2}{\partial x}
+2\rho_{xX}\frac{\partial(g_xg_X)}{\partial X},
\\
A_X
&=
f_X+2\frac{\partial g_X^2}{\partial X}
+2\rho_{xX}\frac{\partial(g_xg_X)}{\partial x},
\end{align}
and
\begin{equation}
\label{eq:A-second}
A_{xx}=g_x^2,
\qquad
A_{xX}=2\rho_{xX}g_xg_X,
\qquad
A_{XX}=g_X^2.
\end{equation}

\section{Coordinate transform and coefficient matching}

Transform from \((x,X)\) to \((y,Y)\), where
\begin{equation}
\label{eq:coordinate-transform}
y=y(x,X),
\qquad
Y=Y(X).
\end{equation}
Since \(Y\) depends only on \(X\),
\[
\frac{dY}{dx}:=Y_x=0,
\qquad
Y_X=\frac{dY}{dX},
\qquad
Y_{XX}=\frac{d^2Y}{dX^2}.
\]
For \(\varphi=\varphi(y,Y,t)\), the chain-rule identities are
\begin{equation}
\label{eq:first-derivatives}
\varphi_x=y_x\varphi_y,
\qquad
\varphi_X=y_X\varphi_y+Y_X\varphi_Y,
\end{equation}
\begin{equation}
\label{eq:second-xx}
\varphi_{xx}=y_x^2\varphi_{yy}+y_{xx}\varphi_y,
\end{equation}
\begin{equation}
\label{eq:second-xX}
\varphi_{xX}
=
y_xy_X\varphi_{yy}
+
y_xY_X\varphi_{yY}
+
y_{xX}\varphi_y,
\end{equation}
and
\begin{align}
\label{eq:second-XX}
\varphi_{XX}
&=
y_X^2\varphi_{yy}
+
2y_XY_X\varphi_{yY}
+
Y_X^2\varphi_{YY}
+
y_{XX}\varphi_y
+
Y_{XX}\varphi_Y.
\end{align}
Substituting \eqref{eq:first-derivatives}--\eqref{eq:second-XX} into \eqref{eq:xX-KFE-expanded} gives
\begin{equation}
\label{eq:yY-KFE-B}
\frac{\partial\varphi}{\partial t}
=
B_0\varphi+B_y\varphi_y+B_Y\varphi_Y+B_{yy}\varphi_{yy}+B_{yY}\varphi_{yY}+B_{YY}\varphi_{YY},
\end{equation}
with
\begin{equation}
\label{eq:B0}
B_0=A_0,
\end{equation}
\begin{equation}
\label{eq:By}
B_y
=
A_xy_x+A_Xy_X+A_{xx}y_{xx}+A_{xX}y_{xX}+A_{XX}y_{XX},
\end{equation}
\begin{equation}
\label{eq:BY}
B_Y=A_XY_X+A_{XX}Y_{XX},
\end{equation}
\begin{equation}
\label{eq:Byy}
B_{yy}
=
A_{xx}y_x^2+A_{xX}y_xy_X+A_{XX}y_X^2,
\end{equation}
\begin{equation}
\label{eq:ByY}
B_{yY}=A_{xX}y_xY_X+2A_{XX}y_XY_X,
\end{equation}
and
\begin{equation}
\label{eq:BYY}
B_{YY}=A_{XX}Y_X^2.
\end{equation}

In the transformed coordinates, the diffusion-side coefficients have the same structural form:
\begin{equation}
\label{eq:yY-diffusion-coefficients}
B_{yy}=g_y^2,
\qquad
B_{yY}=2\rho_{yY}g_yg_Y,
\qquad
B_{YY}=g_Y^2.
\end{equation}
Therefore the diffusion matching equations are
\begin{equation}
\label{eq:diff-match-yy}
g_y^2
=
g_x^2y_x^2
+
2\rho_{xX}g_xg_Xy_xy_X
+
g_X^2y_X^2,
\end{equation}
\begin{equation}
\label{eq:diff-match-yY}
2\rho_{yY}g_yg_Y
=
2\rho_{xX}g_xg_Xy_xY_X
+
2g_X^2y_XY_X,
\end{equation}
and
\begin{equation}
\label{eq:diff-match-YY}
g_Y^2=g_X^2Y_X^2.
\end{equation}
The final equation gives the long-boundary volatility alignment
\begin{equation}
\label{eq:Y-vol-alignment}
Y_X=\frac{g_Y}{g_X},
\end{equation}
where the positive branch preserves the orientation of the long-boundary coordinate. We impose the analogous volatility-aligned orientation for the interior maturity coordinate,
\begin{equation}
\label{eq:y-vol-alignment}
y_x=\frac{g_y}{g_x}.
\end{equation}
Equations \eqref{eq:Y-vol-alignment} and \eqref{eq:y-vol-alignment} are the geometric compatibility conditions of the long-boundary transform. They are the coefficient-matching version of the maturity-loading identity \eqref{eq:maturity-vol-alignment}.

Once both volatility alignment relations are imposed, the two Brownian correlations are not independent. Substituting
\[
Y_X=\frac{g_Y}{g_X},
\qquad
 y_x=\frac{g_y}{g_x}
\]
into the mixed diffusion equation \eqref{eq:diff-match-yY} gives
\[
\rho_{yY}g_yg_Y
=
\rho_{xX}g_xg_X
\left(\frac{g_y}{g_x}\right)
\left(\frac{g_Y}{g_X}\right)
+
g_X^2y_X\left(\frac{g_Y}{g_X}\right).
\]
Therefore
\[
\rho_{yY}g_yg_Y=\rho_{xX}g_yg_Y+g_Xg_Yy_X,
\]
and, away from the degenerate case \(g_Y=0\),
\begin{equation}
\label{eq:yX-forced}
y_X
=
\frac{g_y}{g_X}\left(\rho_{yY}-\rho_{xX}\right).
\end{equation}
Substituting \eqref{eq:y-vol-alignment} and \eqref{eq:yX-forced} into the \(yy\) diffusion equation \eqref{eq:diff-match-yy} yields
\[
g_y^2
=
g_y^2
+
2\rho_{xX}g_y^2\left(\rho_{yY}-\rho_{xX}\right)
+
g_y^2\left(\rho_{yY}-\rho_{xX}\right)^2.
\]
Cancelling the first term on the right hand side and factoring gives
\begin{equation}
\label{eq:rho-branch-constraint}
g_y^2
\left(\rho_{xX}-\rho_{yY}\right)
\left(\rho_{xX}+\rho_{yY}\right)=0.
\end{equation}
Thus, away from the degenerate case \(g_y=0\), the compatible branches are
\begin{equation}
\label{eq:rho-branches}
\rho_{yY}=\rho_{xX}
\qquad\text{or}\qquad
\rho_{yY}=-\rho_{xX}.
\end{equation}
The first branch is the orientation-preserving branch: the coordinate transformation preserves the sign of the Brownian correlation. We therefore use the single structural correlation
\begin{equation}
\label{eq:single-rho-branch}
\rho_{yY}=\rho_{xX}=\rho
\end{equation}
in the admissible affine specification. The opposite branch corresponds to a sign reversal in one coordinate orientation rather than to an independent second correlation parameter.

The \(Y\)-drift coefficient gives the useful separation relation. From \eqref{eq:BY}, using \(Y_X=g_Y/g_X\) and \(A_{XX}=g_X^2\),
\begin{equation}
\label{eq:BY-substitute}
B_Y
=
A_X\frac{g_Y}{g_X}
+
g_X^2\frac{d}{dX}\left(\frac{g_Y}{g_X}\right).
\end{equation}
Substituting \(A_X\) gives
\begin{align}
\label{eq:BY-expanded}
B_Y
&=
\frac{g_Y}{g_X}
\left[
f_X+2\frac{\partial g_X^2}{\partial X}
+2\rho_{xX}\frac{\partial(g_xg_X)}{\partial x}
\right]
+
g_X^2\frac{d}{dX}\left(\frac{g_Y}{g_X}\right).
\end{align}
Since \(g_Y=g_Y(Y)\),
\[
\frac{dg_Y}{dX}
=
\frac{\partial g_Y}{\partial Y}\frac{dY}{dX}
=
\frac{\partial g_Y}{\partial Y}\frac{g_Y}{g_X}.
\]
Therefore
\begin{equation}
\label{eq:quotient-term}
g_X^2\frac{d}{dX}\left(\frac{g_Y}{g_X}\right)
=
g_Y\frac{\partial g_Y}{\partial Y}
-
g_Y\frac{\partial g_X}{\partial X}.
\end{equation}
The quotient rule and the independence assumptions remove the remaining mixed terms, yielding a clean separation between the long-boundary and short-side dynamics. Matching \(B_Y\) against the transformed coefficient
\[
B_Y=f_Y+2\frac{\partial g_Y^2}{\partial Y}+2\rho_{yY}\frac{\partial(g_yg_Y)}{\partial y},
\]
dividing by \(g_Y\), and simplifying gives
\begin{align}
\label{eq:separated-relation-general}
\frac{f_Y}{g_Y}
+
3\frac{\partial g_Y}{\partial Y}
+
\frac{2\rho_{yY}\partial_y(g_yg_Y)}{g_Y}
&=
\frac{f_X}{g_X}
+
3\frac{\partial g_X}{\partial X}
+
\frac{2\rho_{xX}\partial_x(g_xg_X)}{g_X}.
\end{align}
If \(g_Y\) is independent of \(y\), \(g_X\) is independent of \(x\), and the orientation-preserving branch \eqref{eq:single-rho-branch} is used, this reduces to
\begin{equation}
\label{eq:separated-relation-clean}
\frac{f_Y}{g_Y}
+
3\frac{\partial g_Y}{\partial Y}
+
2\rho\frac{\partial g_y}{\partial y}
=
\frac{f_X}{g_X}
+
3\frac{\partial g_X}{\partial X}
+
2\rho\frac{\partial g_x}{\partial x}.
\end{equation}
This is the coefficient-matching separation equation for the long-boundary transform.

\section{Affine admissibility and constant MSNR}

We now impose the affine admissible class. In the original coordinates, let
\begin{equation}
\label{eq:old-affine-g}
g_x(x)=p_x+q_xx,
\qquad
g_X(X)=p_X+q_XX,
\end{equation}
and
\begin{equation}
\label{eq:old-affine-f}
f_x=\gamma g_x,
\qquad
f_X=\delta g_X,
\end{equation}
where \(\gamma\) and \(\delta\) are constants. Thus
\[
q_x=\frac{\partial g_x}{\partial x},
\qquad
q_X=\frac{\partial g_X}{\partial X}.
\] Assume also that the orientation-preserving branch gives the single constant correlation
\[
\rho_{xX}=\rho_{yY}=\rho.
\]
Then
\[
\frac{f_x}{g_x}=\gamma,
\qquad
\frac{f_X}{g_X}=\delta.
\]
Thus the original-coordinate component MSNRs are constant.

In the transformed coordinates, impose the analogous affine structure
\begin{equation}
\label{eq:new-affine-g}
g_y(y)=p_y+q_yy,
\qquad
g_Y(Y)=p_Y+q_YY,
\end{equation}
and
\begin{equation}
\label{eq:new-affine-f}
f_y=\alpha g_y,
\qquad
f_Y=\beta g_Y,
\end{equation}
where \(\alpha\) and \(\beta\) are constants. Thus
\[
q_y=\frac{\partial g_y}{\partial y},
\qquad
q_Y=\frac{\partial g_Y}{\partial Y}.
\] The same \(\rho\) is used in the transformed coordinates by \eqref{eq:single-rho-branch}.
Then
\begin{equation}
\label{eq:component-MSNRs}
\frac{f_y}{g_y}=\alpha,
\qquad
\frac{f_Y}{g_Y}=\beta.
\end{equation}
The transformed-coordinate component MSNRs are therefore constant.

The clean separation relation \eqref{eq:separated-relation-clean} becomes
\begin{equation}
\label{eq:affine-separation}
\beta+3q_Y+2\rho q_y
=
\delta+3q_X+2\rho q_x.
\end{equation}
Equivalently,
\begin{equation}
\label{eq:affine-separation-zero}
(\beta-\delta)+3(q_Y-q_X)+2\rho(q_y-q_x)=0.
\end{equation}
Thus the affine family remains admissible subject to the algebraic compatibility condition \eqref{eq:affine-separation-zero}, the volatility-aligned geometric conditions \eqref{eq:Y-vol-alignment}--\eqref{eq:y-vol-alignment}, and the orientation-preserving correlation branch \eqref{eq:single-rho-branch}. The class is non-unique because many choices of affine coefficients and coordinate maps can satisfy these relations.

The correlated scalar MSNR is the quadratic form determined by the diffusion covariance. In the transformed coordinates, the diffusion matrix without the common factor \(2\) is
\begin{equation}
\label{eq:Sigma-yY}
\Sigma
=
\begin{pmatrix}
g_y^2 & \rho g_yg_Y \\
\rho g_yg_Y & g_Y^2
\end{pmatrix}.
\end{equation}
The squared MSNR is
\begin{equation}
\label{eq:MSNR-quadratic-form}
\mathrm{MSNR}_{y,Y}^2
=
\begin{pmatrix} f_y & f_Y \end{pmatrix}
\Sigma^{-1}
\begin{pmatrix} f_y \\ f_Y \end{pmatrix}.
\end{equation}
Since
\[
\det(\Sigma)=g_y^2g_Y^2(1-\rho^2),
\]
we have
\[
\Sigma^{-1}
=
\frac{1}{g_y^2g_Y^2(1-\rho^2)}
\begin{pmatrix}
g_Y^2 & -\rho g_yg_Y \\
-\rho g_yg_Y & g_y^2
\end{pmatrix}.
\]
Therefore
\begin{align}
\label{eq:MSNR-expanded}
\mathrm{MSNR}_{y,Y}^2
&=
\frac{1}{1-\rho^2}
\left[
\left(\frac{f_y}{g_y}\right)^2
-2\rho\left(\frac{f_y}{g_y}\right)\left(\frac{f_Y}{g_Y}\right)
+
\left(\frac{f_Y}{g_Y}\right)^2
\right].
\end{align}
Using \eqref{eq:component-MSNRs}, this becomes
\begin{equation}
\label{eq:constant-correlated-MSNR}
\boxed{
\mathrm{MSNR}_{y,Y}^2
=
\frac{\alpha^2-2\rho\alpha\beta+\beta^2}{1-\rho^2}.
}
\end{equation}
Hence
\begin{equation}
\label{eq:constant-correlated-MSNR-root}
\boxed{
\mathrm{MSNR}_{y,Y}
=
\left(\frac{\alpha^2-2\rho\alpha\beta+\beta^2}{1-\rho^2}\right)^{1/2}.
}
\end{equation}
Since \(\alpha\), \(\beta\), and \(\rho\) are constants, the scalar correlated MSNR is a global constant.

In the symmetric case \(\alpha=\beta=a\),
\[
\mathrm{MSNR}_{y,Y}^2=\mathrm{MSNR}^2
=
\frac{a^2-2\rho a^2+a^2}{1-\rho^2}
=
\frac{2a^2(1-\rho)}{(1-\rho)(1+\rho)}.
\]
For \(\rho\neq1\),
\begin{equation}
\label{eq:symmetric-MSNR}
\boxed{
\mathrm{MSNR}^2=\frac{2a^2}{1+\rho},
\qquad
\mathrm{MSNR}=|a|\sqrt{\frac{2}{1+\rho}}.
}
\end{equation}
This is the two-boundary analogue of the one-boundary constant MSNR clock.

\section{Robin covariance norm and two-channel MSNR}

The coefficient-matching result concerns admissible geometry. The Robin two-channel implementation gives the corresponding risk loading. In the Robin branch we write the maturity-decay parameter and structural correlation as \(a_R\) and \(\rho_R\), respectively, to allow them to differ from their free-boundary counterparts.

Define the short/interior Robin channel by
\begin{equation}
\label{eq:AR-definition}
A_R(t,\tau)=
x_{t,t}(1+B_t\tau)e^{-a_R^2\tau},
\end{equation}
and define the transported long-boundary channel by
\begin{equation}
\label{eq:CR-definition}
C_R(t,\tau)=
\left(1-e^{-a_R^2\tau}\right)Y_t.
\end{equation}
This is the minimal normalisation of the long-boundary channel.

Let
\begin{equation}
\label{eq:rhoR-correlation}
dW_t^SdW_t^X=\rho_Rdt,
\end{equation}
with \(\rho_R\) constant in the clean Robin admissible baseline. The two-channel diffusion part is
\begin{equation}
\label{eq:two-channel-diffusion}
dx_t(\tau)
=
\sqrt{2}a_RA_R(t,\tau)dW_t^S
+
\sqrt{2}a_RC_R(t,\tau)dW_t^X
+
\cdots.
\end{equation}
The instantaneous variance rate is therefore
\begin{equation}
\label{eq:robin-variance-rate}
v_{\rho_R,R}(t,\tau)=
2a_R^2
\left[
A_R(t,\tau)^2
+
C_R(t,\tau)^2
+
2\rho_RA_R(t,\tau)C_R(t,\tau)
\right].
\end{equation}
Define the Robin covariance norm
\begin{equation}
\label{eq:GrhoRR}
G_{\rho_R,R}(t,\tau)=
\left[
A_R(t,\tau)^2
+
C_R(t,\tau)^2
+
2\rho_RA_R(t,\tau)C_R(t,\tau)
\right]^{1/2}.
\end{equation}
Then
\[
v_{\rho_R,R}(t,\tau)=2a_R^2G_{\rho_R,R}(t,\tau)^2.
\]

The covariance norm is an amplitude. It is not itself an oriented scalar loading unless an additional orientation convention is imposed. For this reason, the drift--variance closure is most cleanly written in squared or absolute form:
\begin{equation}
\label{eq:robin-drift-magnitude}
f_{\rho_R,R}(t,\tau)^2=
a_R^4G_{\rho_R,R}(t,\tau)^2,
\end{equation}
or equivalently
\begin{equation}
\label{eq:robin-drift-absolute}
|f_{\rho_R,R}(t,\tau)|=
a_R^2G_{\rho_R,R}(t,\tau).
\end{equation}
Thus the two-channel covariance geometry determines the admissible drift magnitude associated with the IRFR signal-to-noise clock.

This scalar closure should not be confused with the correlated two-channel MSNR. It is a drift-to-realised-standard-deviation ratio after the two stochastic channels have already been collapsed into the covariance norm. Indeed,
\begin{equation}
\label{eq:robin-scalar-sd-ratio}
\mathcal R_{\rho_R,R}^{\mathrm{sd}}(t,\tau)
:=
\frac{|f_{\rho_R,R}(t,\tau)|}{\sqrt{v_{\rho_R,R}(t,\tau)}}
=
\frac{a_R^2G_{\rho_R,R}(t,\tau)}
{\sqrt{2a_R^2G_{\rho_R,R}(t,\tau)^2}}
=\frac{a_R}{\sqrt{2}},
\end{equation}
whenever \(G_{\rho_R,R}(t,\tau)\neq0\). The constancy of \(\mathcal R_{\rho_R,R}^{\mathrm{sd}}\) is therefore a covariance-norm closure statement, not the full inverse-covariance MSNR calculation.

The genuine Robin two-channel MSNR is obtained before the scalar norm collapse. Define the signed diffusion loadings
\[
g_S(t,\tau)=a_RA_R(t,\tau),
\qquad
g_X(t,\tau)=a_RC_R(t,\tau),
\]
and suppose the corresponding channel drifts satisfy
\begin{equation}
\label{eq:robin-channel-drifts}
f_S=\alpha_R g_S,
\qquad
f_X=\beta_R g_X,
\end{equation}
where \(\alpha_R\) and \(\beta_R\) are constants. The diffusion covariance matrix, after removing the common factor \(2\), is
\[
\Sigma_R
=
\begin{pmatrix}
g_S^2 & \rho_Rg_Sg_X \\
\rho_Rg_Sg_X & g_X^2
\end{pmatrix}.
\]
Thus the squared Robin two-channel MSNR is
\begin{equation}
\label{eq:robin-two-channel-MSNR-form}
\mathrm{MSNR}_{R}^2
=
\begin{pmatrix} f_S & f_X \end{pmatrix}
\Sigma_R^{-1}
\begin{pmatrix} f_S \\ f_X \end{pmatrix}.
\end{equation}
Using \eqref{eq:robin-channel-drifts}, this gives
\begin{equation}
\label{eq:robin-two-channel-MSNR}
\boxed{
\mathrm{MSNR}_{R}^2
=
\frac{
\alpha_R^2-2\rho_R\alpha_R\beta_R+\beta_R^2
}{
1-\rho_R^2
}.
}
\end{equation}
In the symmetric Robin clock case \(\alpha_R=\beta_R=a_R\), and for \(\rho_R\neq1\),
\begin{equation}
\label{eq:robin-symmetric-MSNR}
\boxed{
\mathrm{MSNR}_{R}^2
=
\frac{2a_R^2}{1+\rho_R},
\qquad
\mathrm{MSNR}_{R}
=
|a_R|\sqrt{\frac{2}{1+\rho_R}}.
}
\end{equation}
Thus the genuine Robin two-channel MSNR has the same correlation dependence as the free or Neumann two-channel calculation. The apparently fixed value \(a_R/\sqrt{2}\) in \eqref{eq:robin-scalar-sd-ratio} belongs to the collapsed scalar drift-to-standard-deviation ratio, not to the two-channel inverse-covariance MSNR.

In the perfectly correlated case \(\rho_R=1\),
\[
G_{1,R}(t,\tau)^2=
\left(A_R(t,\tau)+C_R(t,\tau)\right)^2.
\]
Therefore
\[
G_{1,R}(t,\tau)=
\left|A_R(t,\tau)+C_R(t,\tau)\right|.
\]
If the collapsed one-direction loading \(A_R+C_R\) is known to remain non-negative on the region under consideration, then the magnitude identity reduces to the familiar linear expression
\[
|f_{1,R}(t,\tau)|
=a_R^2\left(A_R(t,\tau)+C_R(t,\tau)\right).
\]
More generally, the covariance norm gives the risk amplitude, while the orientation of the drift is inherited from the chosen affine drift convention and boundary orientation. The key invariant statement for the scalar Robin closure is therefore the squared or absolute drift--variance identity, not the introduction of a separate sign branch.

For \(\rho_R\neq1\), the covariance norm is not equal to the linear sum \(A_R+C_R\). The maturity Jacobian remains linear, but the effective risk amplitude is the covariance norm of the two stochastic channels. The correlated MSNR, by contrast, is the inverse-covariance quadratic form \eqref{eq:robin-two-channel-MSNR}.

\section{Asymptotic boundary dynamics and variance floor}

Taking the long-maturity limit of the minimal one-risk-direction boundary law gives
\begin{equation}
\label{eq:long-boundary-sde-X}
dX_t
=
-a^2(X_t-x_\infty)\,dt
+
\sqrt{2}a(X_t-x_\infty)\,dW_t^X.
\end{equation}
Equivalently, for \(Y_t=X_t-x_\infty\),
\begin{equation}
\label{eq:multiplicative-displacement-law}
dY_t
=
-a^2Y_t\,dt
+
\sqrt{2}aY_t\,dW_t^X.
\end{equation}
For \(Y_{t_0}\neq0\),
\begin{equation}
\label{eq:Y-explicit-solution}
Y_t
=
Y_{t_0}
\exp\left(-2a^2(t-t_0)+\sqrt{2}a(W_t^X-W_{t_0}^X)\right).
\end{equation}
The sign of the displacement is preserved. If \(Y_{t_0}=0\), then \(Y_t=0\) thereafter in this minimal multiplicative subclass.

The generator for the minimal long-boundary process is
\begin{equation}
\label{eq:long-boundary-generator}
(\mathcal A\phi)(X)
=
-a^2(X-x_\infty)\phi_X(X)
+
a^2(X-x_\infty)^2\phi_{XX}(X).
\end{equation}
Applying the generator to \(\phi(X)=X\) gives the drift \(-a^2(X-x_\infty)\). Applying it to \((X-X_0)^2\) at \(X=X_0\) identifies the variance coefficient
\[
\sigma_X^2(X)=2a^2(X-x_\infty)^2.
\]
Thus the generator corroborates the asymptotic SDE.

For the Robin implementation, the corresponding long-boundary variance floor is obtained by replacing the generic long-boundary clock \(a\) with the Robin clock \(a_R\). As \(\tau\to\infty\),
\[
A_R(t,\tau)\to0,
\qquad
C_R(t,\tau)\to Y_t.
\]
Therefore the Robin long-end variance floor is
\begin{equation}
\label{eq:long-end-variance-floor}
\lim_{\tau\to\infty}v_{\rho_R,R}(t,\tau)
=
2a_R^2Y_t^2
=
2a_R^2(X_t-x_\infty)^2.
\end{equation}
The cross term does not contribute to the asymptotic floor because the short/interior Robin channel dies out at the long-boundary. Correlation matters at finite maturities, where both channels are active.

\section{Maturity-space Jacobian and stochastic rank}

The completed boundary-field representation separates maturity geometry from state dynamics. In the canonical bridge case,
\begin{equation}
\label{eq:boundary-jacobian}
J(\tau)
=
\left(
 e^{-a^2\tau},
 \tau e^{-a^2\tau},
 1-e^{-a^2\tau}
\right).
\end{equation}
For the state vector
\[
S_t=\left(x_{t,t},\,B_tx_{t,t},\,X_t\right)^\top,
\]
the curve may be written as
\begin{equation}
\label{eq:jacobian-state-representation}
x_t(\tau)=J(\tau)S_t.
\end{equation}
The entries of \(J\) have distinct boundary meanings: short-boundary transport, Robin interior deformation, and long-boundary transport.

The loadings are structural boundary coordinates, not statistical principal components. They are linearly independent. Indeed, if
\[
c_1e^{-a^2\tau}+c_2\tau e^{-a^2\tau}+c_3(1-e^{-a^2\tau})=0
\]
for all \(\tau\ge0\), then letting \(\tau\to\infty\) gives \(c_3=0\). The remaining identity
\[
(c_1+c_2\tau)e^{-a^2\tau}=0
\]
forces \(c_1=c_2=0\).

The Jacobian does not impose a particular stochastic rank. In general,
\[
dS_t=\mu_t\,dt+Q_t\,dW_t,
\]
with state covariance \(\Sigma_t=Q_tQ_t^\top\). The maturity covariance kernel is
\begin{equation}
\label{eq:jacobian-covariance-kernel}
K_t(\tau,\sigma)=J(\tau)\Sigma_tJ(\sigma)^\top.
\end{equation}
A one-direction, two-direction, or higher-rank volatility model corresponds to different restrictions on \(\Sigma_t\). The invariant contribution of the boundary theory is the maturity geometry \(J(\tau)\).

\section{Relation to risk-neutral pricing}

The stochastic long-boundary extension is a physical-field construction. It identifies how a realised long-maturity boundary state can move within the same boundary-field architecture and how drift must be paired with variance when short/interior and long-boundary shocks are not perfectly correlated.

The money-market account is generated by the realised short-end roll \(x_{t,t}\). The stochastic long-boundary \(X_t\) does not directly fund the bank account. Instead, it enters bond prices through the full IRFR strip
\[
\{x_t(u-t):u\in[t,T]\}.
\]
A full no-arbitrage treatment must trace \(X_t\), the Robin interior state, and the boundary-risk covariance structure through the risk-neutral drift restriction and the rolling-maturity transport equation. A companion pricing treatment can impose risk-neutral dynamics on the same volatility-aligned geometry and derive the associated no-arbitrage bond-pricing equations. The present paper does not impose that pricing restriction. It completes the physical boundary-field architecture and identifies the admissible long-boundary geometry.

\section{Empirical Implications}

The coefficient-matching theorem identifies constant correlation as part of the clean affine admissible baseline. Thus the empirical specification should treat \(\rho\) as a structural parameter when testing the admissible model. Rolling or regime-dependent correlations may still be reported as diagnostics, but they should not be confused with the constant-correlation theorem.

The empirical object supplied by the theory is the completed Jacobian
\[
J(\tau)=
\left(
 e^{-a^2\tau},
 \tau e^{-a^2\tau},
 1-e^{-a^2\tau}
\right)
\]
together with the Robin two-boundary covariance norm
\[
G_{\rho_R,R}^2=A_R^2+C_R^2+2\rho_R A_RC_R.
\]
The covariance norm is the correct empirical object for the scalar Robin volatility amplitude. A centred scalar ratio based on the collapsed drift \(a_R^2(A_R+C_R)\) is testing the perfectly correlated Robin special case. If the denominator is estimated from a two-channel covariance structure with \(\rho_R\neq1\), then the corresponding scalar numerator should be interpreted in magnitude as \(a_R^2G_{\rho_R,R}\), not as the signed linear sum \(a_R^2(A_R+C_R)\).

The genuine two-channel MSNR is a different object. It is the inverse-covariance quadratic form and, in the symmetric Robin clock case, depends on correlation according to
\[
\mathrm{MSNR}_{R}^2
=
\frac{2a_R^2}{1+\rho_R}.
\]
Thus empirical tests should distinguish between the scalar covariance-norm drift--standard-deviation ratio and the correlated two-channel Robin MSNR. Instability in a centred MSNR statistic need not refute the boundary geometry. It may instead indicate that realised drift estimation, finite-maturity long-boundary proxies, or time variation in empirical correlation make the direct centred statistic fragile.

A more robust empirical test projects observed curve changes onto the completed Jacobian loadings and estimates the effective covariance-norm loading \(G_{\rho_R,R}\), while separately estimating the Robin structural correlation \(\rho_R\) that enters the two-channel MSNR. The stronger empirical target is whether the completed boundary geometry improves the joint description of shape, volatility, and covariance, with constant \(\rho_R\) used as the Robin structural admissible baseline.

\section{Summary}

The stochastic long-boundary extension separates the equilibrium anchor \(x_\infty\) from the realised long-boundary state \(X_t\). The displacement \(Y_t=X_t-x_\infty\) is transported through maturity by \(1-e^{-a^2\tau}\). Together with the short-boundary loading \(e^{-a^2\tau}\) and the Robin interior loading \(\tau e^{-a^2\tau}\), this gives the completed maturity-space Jacobian
\[
J(\tau)=
\left(
 e^{-a^2\tau},
 \tau e^{-a^2\tau},
 1-e^{-a^2\tau}
\right).
\]

The coefficient-matching calculation shows that the long-boundary transform preserves the affine admissible family when the coordinate geometry is volatility-aligned:
\[
y_x=\frac{g_y}{g_x},
\qquad
Y_X=\frac{g_Y}{g_X}.
\]
Under affine \(g\), constant \(f/g\), and the orientation-preserving constant-correlation branch \(\rho_{yY}=\rho_{xX}=\rho\), the transformed two-boundary system has a global constant MSNR:
\[
\mathrm{MSNR}^2
=
\frac{\alpha^2-2\rho\alpha\beta+\beta^2}{1-\rho^2}.
\]
In the symmetric case \(\alpha=\beta=a\), this reduces to
\[
\mathrm{MSNR}^2=\frac{2a^2}{1+\rho}.
\]

The original exponential maturity-loading intuition is therefore not merely compatible with the long-boundary coefficient matching. It is an example of the required volatility-aligned geometry:
\[
\frac{\partial x_{t,T}}{\partial x_{t,t}}
=
e^{-a^2(T-t)}
=
\frac{g_{t,T}}{g_{t,t}}.
\]
When the long-boundary is stochastic, the observed maturity-\(T\) volatility contains both the short-side loading and the long-boundary loading, combined through their correlation. In the Robin implementation, the covariance norm \(G_{\rho_R,R}\) gives the scalar risk amplitude, while the genuine two-channel MSNR remains the inverse-covariance quadratic form. In the symmetric Robin clock case this gives
\[
\mathrm{MSNR}_{R}^2=\frac{2a_R^2}{1+\rho_R}.
\]

The paper therefore completes the theoretical UBFT/IRFR sequence. The long-boundary does not destroy the non-unique affine admissible class. It closes the geometry by adding the long-boundary coordinate and by requiring volatility-aligned transport through the maturity field.



\begin{thebibliography}{9}

\bibitem{OzaIRFR}
Oza, S. (2026).
\emph{An Axiomatically Consistent Single-Factor Term Structure Model with Admissible Humped Volatility}.
SSRN Working Paper, posted May 2026.
Available at SSRN: \url{https://papers.ssrn.com/sol3/papers.cfm?abstract_id=6454299}.

\bibitem{OzaShortBoundary}
Oza, S. (2026).
\emph{Unified Boundary Field Theory of Interest Rates II: Monetary Policy as a Short-End Boundary Condition}.
Zenodo, DOI {10.5281/zenodo.20963161},
\url {https://doi.org/10.5281/zenodo.20943770},
Working paper.

\bibitem{RitchkenChuang}
Ritchken, P., and Chuang, I. (1999).
Interest rate option pricing with volatility humps.
Review of Derivatives Research, 3, 237–262.

\end{thebibliography}

\end{document}
